\section{Ορισμοί}
\begin{Orismosbox}{Πραγματική Συνάρτηση}
\label{orismos:pragmatikh_synarthsh}
Τι ονομάζουμε πραγματική συνάρτηση με πεδίο ορισμού το σύνολο $A$?
\end{Orismosbox}\apanthsh
Πραγματική συνάρτηση με πεδίο ορισμού ένα σύνολο $ A $ είναι μια διαδικασία (κανόνας) $f$ με την οποία \textbf{κάθε} στοιχείο $ x\in A $ αντιστοιχεί σε \textbf{ένα μόνο} πραγματικό αριθμό $ y\in\mathbb{R} $. Το $ y $ λέγεται \textbf{τιμή} της συνάρτησης $ f $ στο $ x $ και συμβολίζεται $ f(x) $.
\begin{center}
\centering
\begin{tikzpicture}[scale=.6]
\draw(0,0) ellipse (1cm and 1.5cm);
\draw(4,0) ellipse (1cm and 1.5cm);
\draw[fill=\xrwma!30] (4.1,0) ellipse (.6cm and 1.1cm);
\draw[-latex] (0,.2) arc (140:40:2.6);
\tkzDefPoint(0,.2){A}
\tkzDefPoint(4,.2){B}
\tkzDrawPoints(A,B)
\tkzLabelPoint[left](A){{\footnotesize $ x $}}
\tkzLabelPoint[right](B){{\footnotesize $ y $}}
\tkzText(0,1.8){$ A $}
\tkzText(4,1.8){$ B $}
\tkzText(2,1.45){$ f $}
\draw[-latex] (3.5,0) -- (2.7,-1) node[anchor=north east] {\footnotesize $ f\left( A \right)  $};
\end{tikzpicture}
\end{center}
\begin{Orismosbox}{Σύνολο τιμών}
\label{orismos:synolo_timwn}
Τι ονομάζεται σύνολο τιμών μιας πραγματικής συνάρτησης $f$ με πεδίο ορισμού $A$?
\end{Orismosbox}
\apanthsh
Σύνολο τιμών μιας συνάρτησης $ f $ με πεδίο ορισμού $ A $ λέγεται το σύνολο που περιέχει όλες τις τιμές $ f(x) $ της συνάρτησης για κάθε  $ x\in A $. Συμβολίζεται με $ f(A) $ και είναι
\[ f(A)=\{y\in\mathbb{R}:y=f(x)\ \textrm{για κάθε}\ x\in A\} \]
\begin{Orismosbox}{Γραφική παράσταση}
\label{orismos:grafikh_parastash}
Τι ονομάζεται γραφική παράσταση μιας συνάρτησης $f$?
\end{Orismosbox}
\apanthsh
Έστω $f$ μια συνάρτηση με πεδίο ορισμού $A$ και $Oxy$ ένα σύστημα συντεταγμένων στο επίπεδο. Το σύνολο των σημείων $M(x, y)$ για τα οποία ισχύει $y = f(x)$, δηλαδή το σύνολο των σημείων $M(x, f(x)), x\in A$, λέγεται γραφική παράσταση της $f$ και συμβολίζεται συνήθως με $C_f$.
%Γραφική παράσταση μιας συνάρτησης $ f $ με πεδίο ορισμού ένα σύνολο $ A $ ονομάζεται το σύνολο των σημείων της μορφής $ M(x,f(x)) $ για κάθε $ x\in A $. Συμβολίζεται με $ C_f $
\[ C_f=\{M(x,y):y=f(x)\ \textrm{για κάθε}\ x\in A\} \]
\begin{Orismosbox}{Ίσες συναρτήσεις}
\label{orismos:ises_synarthseis}
Πότε δύο συναρτήσεις $f,g$ ονομάζονται ίσες?
\end{Orismosbox}\apanthsh
Δύο συναρτήσεις $ f,g $ λέγονται ίσες όταν
\begin{itemize}
\item έχουν το ίδιο πεδίο ορισμού $ A $ και
\item ισχύει $ f(x)=g(x) $ για κάθε $ x\in A $.
\end{itemize}
Για να δηλώσουμε ότι δύο συναρτήσεις είναι ίσες γράφουμε $ f=g $.
\begin{Orismosbox}{Πράξεις μεταξύ συναρτήσεων}
\label{orismos:praxeis_synarthsewn}
Έστω συναρτήσεις $ f,g $ με πεδία ορισμού $ A,B $ αντίστοιχα. Πως ορίζονται οι συναρτήσεις $f+g, f-g,f\cdot g$ και $\frac{f}{g}$?
\end{Orismosbox}
\apanthsh
Δίνονται δύο συναρτήσεις $ f,g $ με πεδία ορισμού $ A,B $ αντίστοιχα. 
\begin{enumerate}
\item Η συνάρτηση $ f+g $ του αθροίσματος των δύο συναρτήσεων έχει τύπο $ f(x)+g(x) $ και πεδίο ορισμού $ D_{f+g}=A\cap B $.
\item Η συνάρτηση $ f-g $ της διαφοράς των δύο συναρτήσεων έχει τύπο $ f(x)-g(x) $ και πεδίο ορισμού $ D_{f-g}=A\cap B $.
\item Η συνάρτηση $ f\cdot g $ του γινομένου των δύο συναρτήσεων έχει τύπο $ f(x)\cdot g(x) $ και πεδίο ορισμού $ D_{f\cdot g}=A\cap B $.
\item Η συνάρτηση $ \frac{f}{g} $ του πηλίκου των δύο συναρτήσεων έχει τύπο $ \frac{f(x)}{g(x)} $ και πεδίο ορισμού $ D_{\frac{f}{g}}=\{x\in A\cap B:g(x)\neq 0\} $.
\end{enumerate}
Αν $ A\cap B=\varnothing $ τότε οι παραπάνω συναρτήσεις δεν ορίζονται.
\begin{Orismosbox}{Σύνθεση συναρτήσεων}
\label{orismos:synthesh_synarthsewn}
Έστω δύο συναρτήσεις $f, g$ με πεδία ορισμού $A,B$ αντίστοιχα. Τι ονομάζουμε σύνθεση της
$f$ με την $g$?
\end{Orismosbox}
\apanthsh
Αν $f, g$ είναι δύο συναρτήσεις με πεδίο ορισμού $A,B$ αντιστοίχως, τότε ονομάζουμε \bmath{σύνθεση της $f$ με την $g$}, και τη συμβολίζουμε με $g\circ f$, τη συνάρτηση με τύπο
\[ (g\circ f)(x)=g(f(x))\]
Το πεδίο ορισμού της $g\circ f$ αποτελείται από όλα τα στοιχεία $x$ του πεδίου ορισμού της $f$ για τα οποία το $f(x)$ ανήκει στο πεδίο ορισμού της g. Δηλαδή είναι το σύνολο
\[ D_{g\circ f}=\{x\in\mathbb{R}| x\in A\ \textrm{ και }\ f(x)\in B\} \]
Είναι φανερό ότι η $g\circ f$ ορίζεται αν $A_1\neq\varnothing$, δηλαδή αν $f(A)\cap B\neq\varnothing$.
\begin{center}
\begin{tikzpicture}[scale=.6]
\draw(0,0) ellipse (1cm and 1.5cm);
\draw(4,0) ellipse (1cm and 1.5cm);
\begin{scope}
\draw[clip](4,0) ellipse (1cm and 1.5cm);
\draw[fill=\xrwma!30] (5,0) ellipse (1cm and 1.5cm);
\end{scope}
\draw (5,0) ellipse (1cm and 1.5cm);
\draw (9,0) ellipse (1cm and 1.5cm);
\draw (-.1,0)[fill=\xrwma!30] ellipse (.6cm and 1.1cm);
\draw[-latex] (0,.2) arc (140:40:2.95);
\draw[-latex] (4.5,.2) arc (140:40:2.95);
\draw[-latex] (0,.2) arc (140:40:5.9);
\tkzDefPoint(0,.2){A}
\tkzDefPoint(4.5,.15){B}
\tkzDefPoint(9,.15){C}
\tkzDrawPoints(A,B,C)
\tkzLabelPoint[left](A){{\footnotesize $ x $}}
\tkzLabelPoint[below](B){{\footnotesize $ f(x) $}}
\tkzLabelPoint[below](C){{\footnotesize $ g(f(x)) $}}
\tkzText(0,1.8){\footnotesize$ A $}
\tkzText(3.8,1.8){\footnotesize$ f(A) $}
\tkzText(5.2,1.8){\footnotesize$ B $}
\tkzText(2.4,1.55){\footnotesize$ f $}
\tkzText(6.4,1.55){\footnotesize$ g $}
\tkzText(4.5,2.55){\footnotesize$ g\circ f $}
\tkzText(9,1.8){\footnotesize$ g(B) $}
\draw[-latex] (4.5,-0.8) -- (3.2,-1.9) node[anchor=east] {\footnotesize $ f\left( A \right)\cap B  $};
\draw[-latex] (0,-.5) -- (-.8,-1.7) node[anchor=east,xshift=1mm] {\footnotesize $ D_{g\circ f}  $};
\end{tikzpicture}
\end{center}
Για να ορίζεται η συνάρτηση $ g\circ f $ θα πρέπει να ισχύει $ f(A)\cap B\neq\varnothing $.\\\\
(Αντίστοιχα ορίζεται και η σύνθεση $ f\circ g $ με πεδίο ορισμού το $ D_{f\circ g}=\{x\in\mathbb{R}|x\in B\ \ \textrm{και}\ \ g(x)\in A\} $ και τύπο $ (f\circ g)(x)=f(g(x)) $.)\\\\
\begin{Orismosbox}{Γνησίως αύξουσα συνάρτηση}
\label{orismos:gnisiws_auxousa}
Πότε μία συνάρτηση $f$ ονομάζεται γνησίως αύξουσα σε ένα διάστημα του πεδίου ορισμού της?
\end{Orismosbox}\apanthsh
Έστω μια συνάρτηση $ f $ και $ \Delta $ ένα διάστημα του πεδίου ορισμού της. Η $ f $ θα ονομάζεται γνησίως αύξουσα στο $\Delta $ αν για κάθε $ x_1,x_2\in\Delta $ με $ x_1<x_2 $ ισχύει $ f(x_1)<f(x_2) $:
\[ x_1<x_2\Rightarrow f(x_1)<f(x_2) \]
\begin{Orismosbox}{Γνησίως φθίνουσα συνάρτηση}
\label{orismos:gnisiws_fthinousa}
Πότε μία συνάρτηση $f$ ονομάζεται γνησίως φθίνουσα σε ένα διάστημα του πεδίου ορισμού της?
\end{Orismosbox}\apanthsh
Έστω μια συνάρτηση $ f $ και $ \Delta $ ένα διάστημα του πεδίου ορισμού της. Η $ f $ θα ονομάζεται γνησίως φθίνουσα στο $\Delta $ αν για κάθε $ x_1,x_2\in\Delta $ με $ x_1<x_2 $ ισχύει $ f(x_1)>f(x_2) $:
\[ x_1<x_2\Rightarrow f(x_1)>f(x_2) \]
Η $ f $ σε κάθε περίπτωση λέγεται \textbf{γνησίως μονότονη}.
\begin{Orismosbox}{Ολικό μέγιστο}
\label{orismos:oliko_megisto}
Έστω μια συνάρτηση $f$ με πεδίο ορισμού $A$. Τι ονομάζεται ολικό μέγιστο της $f$?
\end{Orismosbox}
\apanthsh
Μια συνάρτηση $ f $ με πεδίο ορισμού ένα σύνολο $ A $ θα λέμε ότι παρουσιάζει ολικό μέγιστο στο $ x_0 $ το $ f(x_0) $ όταν 
\[ f(x)\leq f(x_0)\ \ \textrm{για κάθε}\ \ x\in A \]
\begin{Orismosbox}{Ολικό ελάχιστο}
\label{orismos:oliko_elaxisto}
Έστω μια συνάρτηση $f$ με πεδίο ορισμού $A$. Τι ονομάζεται ολικό ελάχιστο της $f$?
\end{Orismosbox}
\apanthsh
Μια συνάρτηση $ f $ με πεδίο ορισμού ένα σύνολο $ A $ θα λέμε ότι παρουσιάζει ολικό ελάχιστο στο $ x_0 $ το $ f(x_0) $ όταν 
\[ f(x)\geq f(x_0)\ \ \textrm{για κάθε}\ \ x\in A \]
Το ολικό μέγιστο και ολικό ελάχιστο μιας συνάρτησης ονομάζονται \textbf{ολικά ακρότατα}. Το $ x_0 $ λέγεται \textbf{θέση} ακρότατου.
\begin{Orismosbox}{Συνάρτηση $ 1-1 $}
\label{orismos:synarthsh_1_1}
Πότε μια συνάρτηση $f$ με πεδίο ορισμού $A$ ονομάζεται συνάρτηση $1-1$?
\end{Orismosbox}
\apanthsh
Μια συνάρτηση $ f:A\rightarrow\mathbb{R} $ ονομάζεται $ 1-1 $ εάν για κάθε ζεύγος αριθμών $ x_1,x_2\in A $ του πεδίου ορισμού της $ f $ ισχύει \[ x_1\neq x_2\Rightarrow f(x_1)\neq f(x_2) \]
\begin{Orismosbox}{Αντίστροφη συνάρτηση}
\label{orismos:antistrofh_synarthsh}
Έστω $ f:A\to\mathbb{R} $ μια συνάρτηση μία $1-1$ συνάρτηση. Πώς ορίζεται η αντίστροφη συνάρτηση $f^{-1}$ της $f$?
\end{Orismosbox}
\apanthsh
Έστω μια $1-1$ συνάρτηση $ f:A\to\mathbb{R} $ με σύνολο τιμών $ f(A) $. Η συνάρτηση με την οποία κάθε $ y\in f(A) $ αντιστοιχεί σε ένα \textbf{μοναδικό} $ x\in A $ για το οποίο ισχύει $ f(x)=y $, λέγεται αντίστροφη συνάρτηση της $ f $ και συμβολίζεται με $f^{-1}$.
\begin{center}
\begin{tikzpicture}[scale=.6]
\draw(0,0) ellipse (1cm and 1.5cm);
\draw(4,0) ellipse (1cm and 1.5cm);
\draw[fill=\xrwma!50] (4.1,0) ellipse (.6cm and 1.1cm);
\draw[latex-] (0,.2) arc (140:40:2.6);
\tkzDefPoint(0,.2){A}
\tkzDefPoint(4,.2){B}
\tkzDrawPoints(A,B)
\tkzLabelPoint[left](A){{\footnotesize $ x $}}
\tkzLabelPoint[right](B){{\footnotesize $ y $}}
\tkzText(0,1.8){$ A $}
\tkzText(4,1.8){$ B $}
\tkzText(2,1.45){$ f^{-1} $}
\draw[-latex] (3.5,0) -- (2.7,-1) node[anchor=north east] {\footnotesize $ f\left( A \right)  $};
\end{tikzpicture}
\end{center}
\begin{Orismosbox}{Συνεχής συνάρτηση σε σημείο}
\label{orismos:synechhs_se_shmeio}
Πότε μια συνάρτηση $f$ ονομάζεται συνεχής σε ένα σημείο $x_0$ του πεδίου ορισμού της?
\end{Orismosbox}
\apanthsh
Μια συνάρτηση $ f $ ονομάζεται συνεχής σε ένα σημείο $ x_0 $ του πεδίου ορισμού της όταν  \[ \lim_{x\rightarrow x_0}{f(x)}=f(x_0) \]
\begin{Orismosbox}{Συνεχής συνάρτηση}
\label{orismos:synechhs_synarthsh}
Πότε μια συνάρτηση $f$ ονομάζεται συνεχής?
\end{Orismosbox}
\apanthsh
Μια συνάρτηση $ f $ θα λέμε ότι είναι \textbf{συνεχής}, εάν είναι συνεχής σε κάθε σημείο του πεδίου ορισμού της.
\begin{Orismosbox}{Συνεχής συνάρτηση σε διάστημα}
\label{orismos:synechhs_se_diasthma}
Πότε μια συνάρτηση $f$ ονομάζεται συνεχής σε ένα 
\begin{multicols}{2}
\begin{enumerate}
\item ανοικτό διάστημα $(a,\beta)$?
\item κλειστό διάστημα $[a,\beta]$?
\end{enumerate}
\end{multicols}
\end{Orismosbox}
\apanthsh
\vspace{-7mm}
\begin{enumerate}
\item Μια συνάρτηση $ f $ θα λέγεται συνεχής σε ένα \textbf{ανοιχτό} διάστημα $ (a,\beta) $ εάν είναι συνεχής σε κάθε σημείο του διαστήματος.
\item Μια συνάρτηση $ f $ θα λέγεται συνεχής σε ένα \textbf{κλειστό} διάστημα $ [a,\beta] $ εάν είναι συνεχής σε κάθε σημείο του ανοιχτού διαστήματος και επιπλέον ισχύει
\[ \lim_{x\to a^+}{f(x)}=f(a)\ \ \textrm{και}\ \ \lim_{x\to\beta^-}{f(x)}=f(\beta) \]
\end{enumerate}
\begin{Orismosbox}{Εφαπτομένη γραφικής παράστασης}
\label{orismos:efaptomenh_grafikhs_parastashs}
Έστω $f$ μια συνάρτηση και $A(x_0,f(x_0))$ ένα σημείο της $C_f$. Αν υπάρχει το όριο \[\lim_{x\to x_0}{\frac{f(x)-f(x_0)}{x-x_0}}\]
και είναι ένας πραγματικός αριθμός $\lambda$, τότε ορίζουμε ως εφαπτομένη της $C_f$ στο σημείο της $A$, την ευθεία που διέρχεται από το $A$ και έχει συντελεστή διεύθυνσης $\lambda$. Η ευθεία έχει εξίσωση
\[ y-f(x_0)=\lambda(x-x_0) \]
\end{Orismosbox}
\begin{Orismosbox}{Παράγωγος σε σημείο}
\label{orismos:paragwgos_se_shmeio}
Πότε μια συνάρτηση $f$ ονομάζεται παραγωγίσιμη σε ένα σημείο $x_0$ του πεδίου ορισμού της?
\end{Orismosbox}
\apanthsh
Μια συνάρτηση $ f $ λέγεται \textbf{παραγωγίσιμη} σε ένα σημείο $ x_0 $ του πεδίου ορισμού της αν το όριο
\[ \lim_{x\rightarrow x_0}{\frac{f(x)-f(x_0)}{x-x_0}} \]
υπάρχει και είναι πραγματικός αριθμός. Το όριο αυτό ονομάζεται \textbf{παράγωγος} της $ f $ στο $ x_0 $ και συμβολίζεται $ f'(x_0) $.
\begin{Orismosbox}{Παραγωγίσιμη συνάρτηση}
\label{orismos:paragwgisimh_synarthsh}
Πότε μια συνάρτηση $f$, με πεδίο ορισμού $A$, λέγεται παραγωγίσιμη?
\end{Orismosbox}
\apanthsh
Μια συνάρτηση $ f $ θα λέγεται παραγωγίσιμη στο \textbf{πεδίο ορισμού} της ή απλά παραγωγίσιμη, όταν είναι παραγωγίσιμη σε κάθε σημείο $ x_0\in D_f $.
\begin{Orismosbox}{Παραγωγίσιμη συνάρτηση σε διάστημα}
\label{orismos:paragwgisimh_se_diasthma}
Πότε μια συνάρτηση $f$ λέγεται παραγωγίσιμη σε ένα
\begin{multicols}{2}
\begin{enumerate}
\item ανοικτό διάστημα $(a,\beta)$?
\item κλειστό διάστημα $[a,\beta]$?
\end{enumerate}
\end{multicols}
\end{Orismosbox}
\apanthsh
\vspace{-7mm}
\begin{enumerate}
\item Μια συνάρτηση $ f $ θα λέγεται παραγωγίσιμη σε ένα \textbf{ανοικτό} διάστημα $ (a,\beta) $ του πεδίου ορισμού της όταν είναι παραγωγίσιμη σε κάθε σημείο $ x_0\in(a,\beta) $.
\item Μια συνάρτηση $ f $ θα λέγεται παραγωγίσιμη σε ένα \textbf{κλειστό} διάστημα $ [a,\beta] $ του πεδίου ορισμού της όταν είναι παραγωγίσιμη σε κάθε σημείο $ x_0\in(a,\beta) $ και επιπλέον ισχύει
\[ \lim_{x\to a^+}{\frac{f(x)-f(a)}{x-a}}\in\mathbb{R}\ \ \textrm{ και }\ \ \lim_{x\to \beta^-}{\frac{f(x)-f(\beta)}{x-\beta}}\in\mathbb{R} \]
\end{enumerate}
\begin{Orismosbox}{Πρώτη παράγωγος}
\label{orismos:prwth_paragwgos}
Έστω μια συνάρτηση $f$ ορισμένη σε ένα σύνολο $A$. Τι ονομάζεται πρώτη παράγωγος της $f$?
\end{Orismosbox}
\apanthsh
Έστω μια συνάρτηση $ f:Α\to\mathbb{R} $ και έστω $ A_1 $ το σύνολο των σημείων $ x\in A $ για τα οποία η $ f $ είναι παραγωγίσιμη. Η συνάρτηση με την οποία κάθε $ x\in A_1 $ αντιστοιχεί στο $ f'(x) $ ονομάζεται \textbf{πρώτη παράγωγος} της $ f $ ή απλά \textbf{παράγωγος} της $ f $. Συμβολίζεται με $ f' $.
\begin{Orismosbox}{Δεύτερη παράγωγος}
\label{orismos:deyterh_paragwgos}
Έστω μια συνάρτηση $f$ με πεδίο ορισμού $A$. Τι ονομάζεται δεύτερη παράγωγος της $f$. Πως ορίζεται η $\nu-$οστή παράγωγος της $f$?
\end{Orismosbox}
\apanthsh
Έστω $ A_1 $  το σύνολο των σημείων για τα οποία η $ f $ είναι παραγωγίσιμη. Αν υποθέσουμε ότι το $ A_1 $ είναι διάστημα ή ένωση διαστημάτων τότε η παράγωγος της $ f' $, αν υπάρχει, λέγεται δεύτερη παράγωγος της $ f $ και συμβολίζεται με $ f'' $. Επαγωγικά ορίζεται και η $ \nu- $οστή παράγωγος της $ f $ και συμβολίζεται με $ f^{(\nu)} $. Δηλαδή
\[ f^{(\nu)}=\left[f^{(\nu-1)}\right]' \]
\begin{Orismosbox}{Ρυθμός μεταβολής}
\label{orismos:rythmos_metavolhs}
Τι ονομάζουμε ρυθμό μεταβολής ενός ποσού $ y=f(x) $ ως προς ένα ποσό $ x $ σε ένα σημείο $ x_0 $?
\end{Orismosbox}
\apanthsh
Αν δύο μεταβλητά μεγέθη  $ x , y $ συνδέονται με τη σχέση $ y = f(x) $ , όταν  $ f  $ είναι μια συνάρτηση παραγωγίσιμη στο $ x_0 $, τότε ονομάζουμε \textbf{ρυθμό μεταβολής} του $ y $ ως προς το $ x $ στο σημείο  $ x_0 $ την παράγωγο $ f '(x_0) $.
\begin{Orismosbox}{Τοπικό μέγιστο}
\label{orismos:topiko_megisto}
Πότε μια συνάρτηση $f$ με πεδίο ορισμού $A$ παρουσιάζει τοπικό μέγιστο σε ένα σημείο $x_0\in A$?
\end{Orismosbox}
\apanthsh
Μια συνάρτηση $ f $, με πεδίο ορισμού $ A $, θα λέμε ότι παρουσιάζει τοπικό μέγιστο στο $ x_0\in A $ όταν υπάρχει $ \delta>0 $ τέτοιο ώστε
\[ f(x)\leq f(x_0), \ \ \textrm{για κάθε }x\in A\cap(x_0-\delta,x_0+\delta) \]
Το $ x_0 $ λέγεται \textbf{θέση} η σημείο τοπικού μέγιστου, ενώ το $ f(x_0) $ τοπικό μέγιστο της $ f $.
\begin{Orismosbox}{Τοπικό ελάχιστο}
\label{orismos:topiko_elaxisto}
Πότε μια συνάρτηση $f$ με πεδίο ορισμού $A$ παρουσιάζει τοπικό ελάχιστο σε ένα σημείο $x_0\in A$?
\end{Orismosbox}
\apanthsh
Μια συνάρτηση $ f $, με πεδίο ορισμού $ A $, θα λέμε ότι παρουσιάζει τοπικό ελάχιστο στο $ x_0\in A $ όταν υπάρχει $ \delta>0 $ τέτοιο ώστε
\[ f(x)\geq f(x_0), \ \ \textrm{για κάθε }x\in A\cap(x_0-\delta,x_0+\delta) \]
Το $ x_0 $ λέγεται \textbf{θέση} η σημείο τοπικού ελάχιστου, ενώ το $ f(x_0) $ τοπικό ελάχιστο της $ f $.
\begin{Orismosbox}{Τοπικά ακρότατα}
\label{orismos:topika_akrotata}
Τι ονομάζουμε τοπικά ακρότατα μιας συνάρτησης $f$?
\end{Orismosbox}
\apanthsh
Τα τοπικά ελάχιστα και τα τοπικά μέγιστα της $ f $ ονομάζονται τοπικά ακρότατα της $ f $.

\begin{Orismosbox}{Πιθανές θέσεις τοπικών ακροτάτων - Κρίσιμα σημεία}
\label{orismos:krisima_shmeia}
Ποιες είναι οι πιθανές θέσεις τοπικών ακροτάτων μια συνάρτησης $f$? Ποια σημεία ονομάζονται κρίσιμα σημεία?
\end{Orismosbox}
\apanthsh
Οι πιθανές θέσεις τοπικών ακροτάτων μιας συνάρτησης $f$ σε ένα διάστημα $\varDelta$ είναι:
\begin{enumerate}
\item τα εσωτερικά σημεία του $\varDelta$ όπου μηδενίζεται η $f'$
\item τα εσωτερικά σημεία του $\varDelta$ όπου δεν ορίζεται η $f'$
\item τα άκρα του $\varDelta$ αν είναι κλειστά.
\end{enumerate}
Οι θέσεις \textbf{1 και 2} ονομάζονται \textbf{κρίσιμα σημεία} της $f$.

\begin{Orismosbox}{Κυρτή συνάρτηση}
\label{orismos:kyrth_synarthsh}
Έστω μια συνάρτηση $f$ συνεχής σε ένα διάστημα $\Delta$ και παραγωγίσιμη στο εσωτερικό του
$\Delta$ . Πότε λέμε ότι η συνάρτηση $f$ στρέφει τα κοίλα προς τα πάνω ή είναι κυρτή στο $\Delta$?
\end{Orismosbox}
\apanthsh
Μια συνάρτηση $ f $ λέμε ότι στρέφει τα κοίλα προς τα άνω ή είναι κυρτή στο $ \Delta $, αν η $ f' $ είναι γνησίως αύξουσα στο $ \Delta $.
\begin{Orismosbox}{Κοίλη συνάρτηση}
\label{orismos:koilh_synarthsh}
Έστω μια συνάρτηση $f$ συνεχής σε ένα διάστημα $\Delta$ και παραγωγίσιμη στο εσωτερικό του
$\Delta$ . Πότε λέμε ότι η συνάρτηση $f$ στρέφει τα κοίλα προς τα κάτω ή είναι κοίλη στο $\Delta$?
\end{Orismosbox}
\apanthsh
Μια συνάρτηση $ f $ λέμε ότι στρέφει τα κοίλα προς τα κάτω ή είναι κοίλη στο $ \Delta $, αν η $ f' $ είναι γνησίως φθίνουσα στο $ \Delta $.
\begin{Orismosbox}{Σημείο καμπής}
\label{orismos:shmeio_kamphs}
Έστω μια συνάρτηση $f$ παραγωγίσιμη σ΄ ένα διάστημα $(a,\beta)$, με εξαίρεση ίσως ένα σημείο
του $x_0$. Πότε το σημείο $A(x_0,f(x_0))$ λέγεται σημείο καμπής της γραφικής παράστασης της
$f$?
\end{Orismosbox}
\apanthsh
Έστω μια συνάρτηση $ f $ παραγωγίσιμη σε ένα διάστημα $ (a,\beta) $, με εξαίρεση ίσως ένα σημείο $ x_0 $. Αν:
\begin{itemize}[itemsep=0mm]
\item η $ f $ είναι κυρτή στο $ (a,x_0) $ και κοίλη στο $ (x_0,\beta) $ η αντιστρόφως και
\item η $ C_f $ έχει εφαπτομένη στο $ A(x_0,f(x_0)) $
\end{itemize}
τότε το σημείο $ A(x_0,f(x_0)) $ λέγεται \textbf{σημείο καμπής} της $ C_f $.
\begin{Orismosbox}{Πιθανές θέσεις σημείων καμπής}
\label{orismos:pithanes_theseis_shmeiwn_kamphs}
Ποιες είναι οι πιθανές θέσεις σημείων καμπής μια παραγωγίσιμης συνάρτησης $f$?
\end{Orismosbox}
\apanthsh
Οι πιθανές θέσεις σημείων καμπής μιας συνάρτησης $f$ σε ένα διάστημα $\varDelta$ είναι:
\begin{enumerate}
\item τα εσωτερικά σημεία του $\varDelta$ όπου μηδενίζεται η $f''$
\item τα εσωτερικά σημεία του $\varDelta$ όπου δεν ορίζεται η $f''$.
\end{enumerate}
\begin{Orismosbox}{Κατακόρυφη ασύμπτωτη}
\label{orismos:katakoryfh_asymptwth}
Πότε η ευθεία $x=x_0$ ονομάζεται κατακόρυφη ασύμπτωτη της γραφικής παράστασης μιας συνάρτησης $f$?
\end{Orismosbox}
\apanthsh
Αν ένα τουλάχιστον από τα όρια $ \lim\limits_{x\to x_0^-}{f(x)},\lim\limits_{x\to x_0^+}{f(x)} $ ισούται με $ \pm\infty $ τότε η ευθεία $ x=x_0 $ λέγεται \textbf{κατακόρυφη ασύμπτωτη} της $ C_f $.
\begin{Orismosbox}{Οριζόντια ασύμπτωτη}
\label{orismos:orizontia_asymptwth}
Πότε η ευθεία $y=y_0$ ονομάζεται οριζόντια ασύμπτωτη της γραφικής παράστασης μιας συνάρτησης $f$?
\end{Orismosbox}
\apanthsh
Αν $ \lim\limits_{x\to +\infty}{f(x)}=l $ (αντιστοίχως $ \lim\limits_{x\to -\infty}{f(x)}=l $) τότε η ευθεία $ y=l $ λέγεται \textbf{οριζόντια ασύμπτωτη} της $ C_f $ στο $ +\infty $ (αντίστοιχα στο $ -\infty $).
\begin{Orismosbox}{Πλάγια ασύμπτωτη}
\label{orismos:plagia_asymptwth}
Πότε η ευθεία $y=\lambda x+\beta$ ονομάζεται ασύμπτωτη της γραφικής παράστασης μιας συνάρτησης $f$?
\end{Orismosbox}
\apanthsh
Η ευθεία $ y=\lambda x+\beta $ λέγεται ασύμπτωτη της $ C_f $ στο $ +\infty $ (αντιστοίχως στο $ -\infty $) αν και μόνο αν
\[ \lim\limits_{x\to +\infty}{[f(x)-(\lambda x+\beta)]=0} \]
αντίστοιχα στο $ -\infty $ αν 
\[ \lim_{x\to -\infty}{[f(x)-(\lambda x+\beta)]=0} \]
\begin{itemize}[itemsep=0mm]
\item Αν $ \lambda=0 $ η ασύμπτωτη είναι οριζόντια.
\item Αν $ \lambda\neq 0 $ η ασύμπτωτη είναι πλάγια.
\end{itemize}
\begin{Orismosbox}{Αρχική συνάρτηση}
\label{orismos:arxikh_synarthsh}
Έστω μια συνεχής συνάρτηση $f$ ορισμένη σε ένα διάστημα $\Delta$. Τι ονομάζεται αρχική συνάρτηση ή παράγουσα της $f$ στο $\Delta$?
\end{Orismosbox}
\apanthsh
Αρχική συνάρτηση ή παράγουσα μιας συνεχούς συνάρτησης $f$ σε ένα διάστημα $\varDelta$ του πεδίου ορισμού της, ονομάζεται κάθε παραγωγίσιμη συνάρτηση $F$ για την οποία ισχύει
\[ F'(x)=f(x)\ \ ,\ \ \text{για κάθε }x\in\varDelta \]
